\section{A distributional local law on the finite graph}
\label{sec:local-law}

We now transfer the tree recursion to the finite random graph.  The
quantity to be approximated is an empirical distribution of Green
functions: the diagonal at a typical vertex retains randomness from its
nearby potentials even when the graph is large.  The main step is to
show that the empirical distribution nearly satisfies the independent
branch recursion.  Independence will be obtained by sampling and
switching edges, rather than imposed on the branches of the finite graph.
Throughout this section, $w>0$ is fixed below the threshold in
Proposition~\ref{prop:tree-stability}; constants may depend on $w$.
We write
\[
 \zeta=e+i\eta,\qquad e\in J,\qquad 0<\eta\le h,
 \qquad y=n\eta.
\]
A positive exponent denoted by $c$ may be decreased in successive estimates.
All such exponents are independent of $n,e,\eta$ in this range.

For an oriented edge $(v,v')$, let $H^{[vv']}$ be the Hamiltonian obtained
by deleting that edge, keeping every potential, and let
$G^{[vv']}(\zeta)=(H^{[vv']}-\zeta)^{-1}$.  Define the random probability
measures on $\DD$ by
\begin{equation}\label{eq:empirical-cavity-laws}
 P_H=\frac1{nd}\sum_{(v,v')}\delta_{X(G^{[vv']}_{vv}(\zeta))},
 \qquad
 P_H^0=\frac1n\sum_v\delta_{X(G_{vv}(\zeta))}.
\end{equation}
The first sum is over all oriented edges, and $X$ uses $q(e)$ as in
Section~\ref{sec:tree}.  Let $Q_\zeta=\Law(Y(\zeta))$ be the tree cavity
law in disk coordinates, and let
$Q_\zeta^0=\Law(X(m_o(\zeta)))$ be the full-root law.

\begin{proposition}[Distributional local law]\label{prop:local-law}
There are $c,C>0$ and $y_0<\infty$ such that, for $y=n\eta\ge y_0$,
\begin{equation}\label{eq:distributional-local-law}
 \PP\left\{
 W_2(P_H,Q_\zeta)+W_2(P_H^0,Q_\zeta^0)>y^{-c}
 \right\}\le Cy^{-c}.
\end{equation}
Here $W_2$ is transportation distance for the Euclidean metric on the
disk.  The estimate is uniform in $e\in J$ and $0<\eta\le h$.
\end{proposition}

We prove this proposition after developing the sampling estimates and
the empirical recursion.  The estimates are intentionally only powers
of $n\eta$; later we will first pass to a microscopic limit and then
let its imaginary argument grow.

\subsection{Averaging potentials and sampling resolvent entries}

The following bounds hold without a local law.  Their uniformity under
bounded graph modifications is what permits their use during switching.

\begin{lemma}[Spectral averaging and compression tails]
\label{lem:spectral-averaging}
Let $H=A+\diag(V_v)$, where $A$ is a real symmetric matrix independent
of the iid uniform potentials $V_v\in[-w,w]$.  For $\Im\zeta>0$,
\begin{equation}\label{eq:site-averaging}
 \EE\,\Im G_{vv}(\zeta)\le C_w,
 \qquad
 \EE\,\Tr\mathbf1_I(H)\le C_w n|I|
\end{equation}
for every interval $I\subset\RR$.  If $S$ consists of at most $k$ vertices
chosen without reference to the potentials, then
\begin{equation}\label{eq:compression-tail}
 \PP\{\|G_{SS}(\zeta)\|>D\}\le C_{w,k}D^{-1/2},\qquad D\ge1.
\end{equation}
The same conclusions hold after any graph modifications chosen
independently of the potentials.  The compression estimate also holds
for the inverse of a fixed strictly dissipative $k\times k$ matrix plus
an iid uniform real diagonal; its constant is independent of that matrix.
\end{lemma}

\begin{proof}
Condition on $A$ and all potentials other than $V_v$.  Schur complementation
gives
\begin{equation}\label{eq:intercept}
 G_{vv}(\zeta)=(V_v-A_v(\zeta)-iB_v(\zeta))^{-1},
 \qquad B_v(\zeta)>0,
\end{equation}
where $A_v,B_v$ do not depend on $V_v$.  Integration of the Poisson
kernel over the allowed interval gives
$\EE_{V_v}\Im G_{vv}\le\pi/(2w)$.  Integrating in the real spectral
parameter, followed by Stieltjes inversion, gives the second inequality
in \eqref{eq:site-averaging}.  This is the usual rank-one spectral-averaging
proof of the Wegner estimate \cite{Wegner1981}; the arbitrary finite-matrix
scope is also recorded in \cite[Lemma~3.2]{AizenmanWarzel2009}.

For the last assertion, write the inverse as
$(T+\diag(v_1,\ldots,v_k))^{-1}$ with $\Im T<0$.  Use $v_1$ and the
differences $v_j-v_1$ as coordinates.  The joint density is bounded,
all the differences range over a fixed compact set, and the remaining
variable adds a scalar multiple of the identity on a bounded interval.
For any unit vector $a$, the function
\[
 f(z)=\langle a,(T_0-z)^{-1}a\rangle,\qquad z\in\HH,
\]
where $\Im T_0<0$, takes values in $\HH$ and satisfies $|f(i)|\le1$.
The function
$u(z)=\Re(e^{-i\pi/4}\sqrt{f(z)})$ is positive harmonic and satisfies
$u(z)\ge2^{-1/2}|f(z)|^{1/2}$.  The Poisson inequality at $i$, whose
kernel is bounded below on each fixed bounded real interval, therefore
bounds the integral of $|f(t)|^{1/2}$ on that interval.  This remains
valid by boundary limits if necessary.  Polarization, a fixed basis,
and equivalence of norms in dimension $k$ give a uniform bound on the
integral of the half power of the inverse norm.  Markov's inequality
proves \eqref{eq:compression-tail}.  To apply this to $G_{SS}$, condition
also on the potentials outside $S$ and take its Schur complement.  Its
effective matrix has strictly negative imaginary part.  All conditioning
and integration arguments remain valid for graph-selected sets and for
graph modifications independent of the potentials.
\end{proof}

\begin{lemma}[Ward estimate for independently sampled locations]
\label{lem:ward}
Let $v,v'$ be independent uniform vertices, conditionally on $H$.  Then
\begin{equation}\label{eq:sampled-ward}
 \EE_{H,v,v'}|G_{vv'}(\zeta)|^2
 =\frac1{n\eta}\EE\,\Im\left(\frac1n\Tr G(\zeta)\right)
 \le\frac{C_w}{y}.
\end{equation}
The same estimate applies to vertices read at prescribed positions of
two independently sampled oriented edges or finite ordered
nonbacktracking tree explorations.
\end{lemma}

\begin{proof}
The identity $G^*G=\eta^{-1}\Im G$, followed by averaging its diagonal,
gives the equality.  Lemma~\ref{lem:spectral-averaging} gives the inequality.
For the final assertion, begin each exploration at a uniform vertex or
uniform oriented edge.  At each occurrence of a vertex in the
nonbacktracking lift, independently order the allowed outgoing stubs,
excluding the incoming stub.  In an edge-rooted exploration, exclude the
marked base edge at its initial vertex as well.  Every fixed position
then has uniform vertex marginal, because the graph is regular.
Different explorations are independent conditional on the graph.
This remains true when their images intersect or fail to form trees;
no claim of independence among positions within one exploration is used.
\end{proof}

In a bounded collection of these samples, a collision or a failure of
any prescribed fixed-radius tree neighborhood has probability $O(n^{-1})$.
One can expose the configuration matching one pair at a time: each
exposure has probability $O(n^{-1})$ of hitting the bounded set already
seen.  Conditioning on simplicity multiplies the bound by a fixed
constant, since its positive limiting probability for fixed $d$ follows
from Bollob\'as's enumeration formula \cite{Bollobas1980}; see also
\cite[Lemma~2.1]{bu}.
Ordering seeds may be retained as marks in this exploration.

\subsection{Switching the boundary of a finite tree}

Local edge switchings are also used to recover probabilistic averaging
in local laws for random regular graphs
\cite{BKY2017Local,BHY2019Local}.  Here they will compare a random
empirical cavity law with its own independent-input recursion; the
potential-dependent target is the distributional tree law.

For a probability measure $P$ on $\DD$, let $\Phi_j(P)$ be the law
obtained from a forward tree with $j$ levels: assign iid potentials to
its internal vertices, assign iid $P$-distributed disk values at its
$b^j$ leaves independently of those potentials, and apply $F_j$.
Thus internal vertices are at distances $0,\ldots,j-1$ from the root
and the input leaves are at distance $j$.  Define $\Phi_j^0(P)$ with
$d$ children at the root and $b$ at all later internal vertices.
We also use the augmented laws which record, in their prescribed order,
the potentials at all internal vertices together with the output disk
value.  Empirical observations record the same coordinates from the
ordered finite-graph exploration; on failure to form a tree they are
assigned a fixed default value.  These added coordinates have their
ordinary Euclidean metric.

\begin{proposition}[Empirical recursion]\label{prop:empirical-recursion}
For each fixed $j\ge1$ there are $c,C,y_0>0$ such that, when $y\ge y_0$,
\begin{equation}\label{eq:empirical-recursion}
 \PP\{W_2(P_H,\Phi_j(P_H))>y^{-c}\}\le Cy^{-c}.
\end{equation}
The same assertion holds with $P_H^0$ on the left and
$\Phi_j^0(P_H)$ on the right.  Both assertions also hold for the
augmented observations just defined.  Probabilities in these estimates
average over the graph and potentials; the measures being compared
are conditional laws obtained by sampling on the realized $H$.
\end{proposition}

Fix one of the cavity, full-root, or augmented observations in the
proposition.  For a real test $f$ on its coordinate space, bounded by one
and with Lipschitz constant at most one, write $a_H$ for its expectation under that
empirical observation law and $d_H$ for its expectation under the
corresponding recursion law.  Our immediate objective is
\[
                         \EE|a_H-d_H|^2\le Cy^{-c}.
\]
This compares the two laws on the same realized graph; comparison of
their ensemble averages would not suffice for tree stability.

The marked switch below preserves the joint law of the graph and its
sampled observations.  It couples one or two observations in a common
switched graph $\widehat H$ with recursion outputs that are independent
conditional on the original graph $H$.  Two observations will therefore compare
$\EE a_H^2$ with $\EE d_H^2$.  For the mixed moment
$\EE(a_Hd_H)$ we also need to control $d_{\widehat H}-d_H$:
the prediction on the original graph must be compared with the
prediction on the switched graph.

\paragraph{The marked switch.}
Represent the simple graph by a configuration matching conditioned on
simplicity, with $d$ distinguishable stubs at every vertex.  Choose a
base vertex, or an outgoing base stub whose matched edge will be deleted
for the cavity observation.  Explore its internal tree of $j$ levels,
using the independent ordering seeds from Lemma~\ref{lem:ward}.
Each pending downward edge can be written $p-t$, where $p$ is internal
and $t$ is at the input level.  For each pending stub independently choose
a uniform outgoing far stub, with matched edge denoted $u-v$, and replace
\begin{equation}\label{eq:boundary-switch}
 (p-t,\ u-v)\quad\longmapsto\quad(p-u,\ t-v).
\end{equation}
The operation is performed on the indicated stubs, so several pending
edges at one vertex are distinguished.  Keep the base mark, internal
connecting stubs, and ordering seeds.  Replace the far outgoing mark at
$u$ by the stub at $t$ that was formerly paired with $p$.  On reversal,
this new far mark specifies $t-v$, and the pending stub at $p$ now
specifies $p-u$; a second application restores the old pairs and marks.
Figure~\ref{fig:local-switch} distinguishes this switch from a further
re-pairing used below.

\begin{figure}[t]
\centering
\begin{tikzpicture}[scale=.92,vertex/.style={circle,fill,inner sep=1.5pt},
                    every node/.style={font=\small},>=Stealth]
 \node at (1.4,1.7) {boundary switch};
 \node[vertex,label=left:$p$] (p1) at (0,1) {};
 \node[vertex,label=right:$t$] (t1) at (2,1) {};
 \node[vertex,label=left:$u$] (u1) at (0,0) {};
 \node[vertex,label=right:$v$] (v1) at (2,0) {};
 \draw[thick] (p1)--(t1) (u1)--(v1);
 \node[circle,draw=magenta!70!black,fill=white,line width=1pt,
       inner sep=2pt] at (.4,0) {};
 \draw[->] (2.5,.5)--(3.3,.5);
 \node[vertex,label=left:$p$] (p2) at (3.8,1) {};
 \node[vertex,label=right:$t$] (t2) at (5.8,1) {};
 \node[vertex,label=left:$u$] (u2) at (3.8,0) {};
 \node[vertex,label=right:$v$] (v2) at (5.8,0) {};
 \draw[thick,blue!65!black] (p2)--(u2) (t2)--(v2);
 \node[circle,draw=magenta!70!black,fill=white,line width=1pt,
       inner sep=2pt] at (5.8,.75) {};
 \node[align=center] at (4.8,-.65) {far mark moves\ from $u$ to $t$};
 \begin{scope}[yshift=-3.3cm]
 \node at (1.4,1.7) {re-pairing before two cuts};
 \node[vertex,label=left:$u$] (u3) at (0,1) {};
 \node[vertex,label=right:$v$] (v3) at (2,1) {};
 \node[vertex,label=left:$a$] (a3) at (0,0) {};
 \node[vertex,label=right:$b$] (b3) at (2,0) {};
 \draw[thick,dashed] (u3)--(v3) (a3)--(b3);
 \draw[->] (2.5,.5)--(3.3,.5);
 \node[vertex,label=left:$u$] (u4) at (3.8,1) {};
 \node[vertex,label=right:$v$] (v4) at (5.8,1) {};
 \node[vertex,label=left:$a$] (a4) at (3.8,0) {};
 \node[vertex,label=right:$b$] (b4) at (5.8,0) {};
 \draw[thick,dashed,blue!65!black] (u4)--(b4) (a4)--(v4);
 \node[align=center] at (4.8,-.65) {delete both dashed edges:\ the remaining graph is identical};
 \end{scope}
\end{tikzpicture}
\caption{The upper switch attaches a sampled cavity at $u$ to the
internal vertex $p$, retaining a reversible stub mark.  The open circles
mark the selected outgoing stubs at $u$ before the switch and $t$ after it.
The lower re-pairing is used only to estimate Green entries between the two
endpoints of a cut edge.  After both displayed edges are removed the
graph is unchanged, while $u$ and $v$ belong to different sampled edge
groups.  The drawings specify stub pairings; they assert no independence
of neighboring Green functions.}
\label{fig:local-switch}
\end{figure}

Apply the switch only when the internal exploration has no collisions,
the required parent and pending stub data are distinct except for shared
internal vertices, the far pairs have distinct data disjoint from the
explored data, all replacements are simple, and the inverse instructions
are applicable and restore the input.  On the complementary set define
the transformation to be the identity.  This definition makes it an
involution of the marked sample space.  Its nontrivial part preserves
weights: the underlying matching, each base mark, each far outgoing
stub, and each ordering seed have the same weight before and after the
operation.  A bounded collection of initially separated tree
neighborhoods satisfies all the requirements.  In such neighborhoods no
new edge duplicates an old edge, and reversal reads the same pending
positions and distinct far data.  Consequently the excluded set has
probability $O(n^{-1})$.

This also applies to two independently sampled bases, whose internal
trees are disjoint outside an $O(n^{-1})$ event.  The two observations
use one common switched graph; a cavity observation deletes only its
own marked base edge from that graph, not the other observation's edge.  We do not require
additional acyclicity among exterior vertices as a condition of
reversal.  Thus bounded switch averages may ignore failures at cost
$O(n^{-1})$ and otherwise use the original joint law of the graph and
its marks.  Potentials are attached to vertices and are never used in
choosing a switch, so their iid law is preserved.

\paragraph{Resolvent comparison after switching.}
For the next argument, say that a scalar or fixed-dimensional matrix
error is small if it has norm at most $y^{-a}$ outside an event of
probability at most $Cy^{-a}$, for some $a>0$.  There are only boundedly
many entries and operations, with bounds depending on $j,d$.
Lemma~\ref{lem:spectral-averaging} permits multiplication of a small
error by every fixed number of Green entries that occurs below.
Indeed, if the original exponent is $a$ and at most $N$ factors occur,
restrict each factor to at most $y^\epsilon$ with
$\epsilon<a/(2N)$.  The product still has a negative-power bound, and
the union of tail exceptions has a negative-power probability.
Also $n^{-a}$ is small in this sense since $y\le hn$.

Partition the selected coordinates according to their independent
samplings: one group for each base and its internal vertices, original
boundary endpoints, and designated parent if present; one group for
each far pair.  Lemma~\ref{lem:ward} and Markov's inequality show that
in the unmodified Green matrix all entries joining different groups
are small.  First sever every pending edge and every far edge, and
also sever the base edge for a cavity observation.  We claim that the
cross-group entries stay small and that each far-pair block stays close
to its block with only its own edge deleted.

Here is the finite-rank calculation behind this claim.  Suppose a
bounded update $B$ is supported in a single group $S$, let
$Q=G_{SS}$, and let $Q'=G'_{SS}$ be its updated compression.  The
identities
\begin{equation}\label{eq:group-update}
 \begin{split}
 (I+BQ)^{-1}&=I-BQ',\\
 G'_{UW}-G_{UW}&=-G_{US}(I+BQ)^{-1}B G_{SW}
 \end{split}
\end{equation}
hold for any two coordinate sets $U,W$.  Both $Q$ and $Q'$ have the
bounds allowed above, since the updates and coordinate choices are
independent of the potentials.  If $U$ and $W$ are distinct groups,
at least one factor crossing to $S$ is small; if both belong to a
different group, both crossing factors are small.  Formula
\eqref{eq:group-update} also applies when one of $U,W$ is $S$.
Severing one group at a time proves the claim.  To compare a given
far block with its individual single-cut block, perform that group's
cut first.

We still need smallness between the two ends of a severed far edge.
This does not follow from the Ward estimate for independent locations,
since the endpoints of an uncut edge were sampled together.  Introduce
another independently sampled far pair, if one is not already available,
and write the two pairs as $(u,v)$ and $(a,b)$.  Re-pair them as
$(u,b)$ and $(a,v)$, and then sever both pairs.  The resulting double-cut
graph is exactly the same as when $(u,v)$ and $(a,b)$ were severed.
The re-pairing is an equal-weight marked switch with $O(n^{-1})$
failure probability, by the preceding construction.  In its new marked
representation, $u$ and $v$ belong to different edge groups.  The
cross-group estimate already proved therefore makes their entry in
the double-cut graph small.  If the second pair was introduced solely
for this argument, restore its original edge: the cross entries
between the original groups $\{u,v\}$ and $\{a,b\}$ are small, so \eqref{eq:group-update} preserves the
conclusion.  This proves the required estimate without any assumption
about a cut edge's two endpoints.

At this point the internal tree to be observed is isolated.  Restore
all exterior replacement edges $t-v$, postponing attachment of this
internal tree until the end.  For two bases, also perform the other
base's updates; its parent edge is restored to its ordinary uncut
status because the present cavity observation deletes only its own
base edge from the common switched graph.  Every cross entry from a new
leaf $u$ to these update coordinates is small, including the entry
to its former partner $v$ just treated by re-pairing.  Here take $S$ to be exactly the coordinates changed in this exterior
step: the relevant $t$'s and $v$'s and, when present, the other base's
update coordinates.  This set excludes the present input leaves and
need not be a union of original groups.  Formula
\eqref{eq:group-update} applies to this support as well: every entry
from a present leaf to $S$ is small, and the inverse factor on $S$
still has the fixed-dimensional compression bound.  Therefore the
boundary Green matrix on the new leaves differs by a small error from
the diagonal matrix whose entries are
\[
             G^{[uv]}_{uu}(\zeta)
\]
in the original graph with the individual far edge deleted.

Finally add the $p-u$ edges.  On the internal tree the exact Schur
matrix uses this boundary compression.  The inverse of that Schur
matrix is the internal compression of the final graph resolvent,
so its norm is at most any fixed small positive power of $y$ outside
a negative-power exception, by Lemma~\ref{lem:spectral-averaging}.
Replacing the boundary matrix by its diagonal approximation changes
the Schur matrix by a small error.  The inverse identity, or its
Neumann expansion once the product of inverse norm and error is less
than $1/2$, shows that its inverse changes by a small error as well.
With the diagonal substitution, successive leaf elimination is exactly
the recursion $F_j$, or its full-root version, with the original internal
potentials.  The map $X$ has bounded derivative on $\HH$, uniformly
for $e\in J$, and so conversion to disk coordinates preserves the
error estimate.  This establishes the claimed recursion comparison
for each of one or two switched observations on the same switched data.

\paragraph{Replacing empirical internal potentials by iid potentials.}
The preceding comparison still uses potentials at an empirically
sampled base.  We next replace their empirical law by the deterministic
iid law while keeping independence from the sampled cavity inputs.
Let $k_j$ be the number of internal vertices.  Bin each potential
coordinate on a mesh $n^{-\alpha}$, with $\alpha>0$ sufficiently small
for this fixed $k_j$.  Let $\mu_H$ be the empirical law of the internal
potential vector, including the ordering randomness and a default on
tree failures, and let $\mu$ be the iid potential law.
For each product bin $B$, one- and two-base explorations give
\[
 |\EE\mu_H(B)-\mu(B)|=O(n^{-1}),\qquad
 |\EE\mu_H(B)^2-\mu(B)^2|=O(n^{-1}).
\]
Indeed disjoint valid neighborhoods have independent potential vectors,
and all failures cost $O(n^{-1})$.  Hence
$\EE|\mu_H(B)-\mu(B)|=O(n^{-1/2})$.
There are $O(n^{\alpha k_j})$ bins, so choosing
$\alpha k_j<1/4$, for example, bounds the expected total variation
of the binned laws by a negative power of $n$.

Couple the bins maximally and, on a matched bin, fill in an iid vector
with its conditional distribution inside that bin.  Thus internal
potential vectors can be coupled with coordinate error $O(n^{-\alpha})$
outside an event of negative-power probability.  Given $H$, this
coupling uses only the base sampling and additional randomness.  It
is consequently independent of all far-edge samples, which give iid
single-cut values with law $P_H$ conditional on $H$.  For two bases
use independent such couplings conditional on $H$; their target
potential vectors are iid and independent of both collections of
inputs.  Dependence of $P_H$ on the potentials creates no difficulty
in this conditional assertion.

Changing the finitely many internal potentials by the coupled error
changes the finite-tree recursion by a small error.  This follows
from the same Schur inverse perturbation just used.  Alternatively,
with iid internal potentials the effective finite-tree inverse
satisfies the final assertion of Lemma~\ref{lem:spectral-averaging},
uniformly in its dissipative leaf data.  The comparison also records
the internal potential coordinates themselves, which proves the
augmented version of the approximation.

We have thus constructed the coupling announced above.  For one or
two bases, $\widehat O_i$ is the observation at base $i$ in the common
switched graph, and $R_i$ is the recursion output formed from its
original far-edge samples and the coupled iid internal potentials.
The bounded test $f$ turns the comparison errors into
\begin{equation}\label{eq:switched-observation-comparison}
 \EE|f(\widehat O_i)-f(R_i)|\le Cy^{-c}.
\end{equation}
Conditional on $H$, the $R_i$ are independent and
$\EE[f(R_i)\mid H]=d_H$.  Marked-switch invariance preserves the
joint law of $\widehat H$ and its sampled observations
$\widehat O_i$, while their comparison partners $R_i$ use the original $H$.

\paragraph{Stability of the recursion prediction under a switch.}
For the mixed-moment comparison, we need
\begin{equation}\label{eq:multiplier-stability}
                       \EE|d_{\widehat H}-d_H|\le Cy^{-c}.
\end{equation}
To prove it, use fresh independent far-edge samples for the inputs
defining $d_H$.  Couple the samples in $H$ and $\widehat H$ by their
outgoing stubs; they agree away from the bounded set of modified stubs,
and overlap with update data has probability $O(n^{-1})$.
The old Green entries between each fresh sample group and the update
groups are small by Lemma~\ref{lem:ward}.  Sever its own edge first
and apply \eqref{eq:group-update} through the remaining cuts and
updates, using the union of their coordinates as the support $S$
when an update joins different original groups.  These cross entries stay small, and its single-cut diagonal
changes by a small error, in half-plane as well as disk coordinates.
Use the same fresh iid internal potentials for the two predictions.
The dissipative finite-tree compression bound in
Lemma~\ref{lem:spectral-averaging} and inverse perturbation compare
the two outputs.  Boundedness of $f$ turns the exceptional probabilities
and the good-event errors into \eqref{eq:multiplier-stability}.

\paragraph{Two observations and transportation distance.}
We can now complete the proof of Proposition~\ref{prop:empirical-recursion}.
With two bases, switch invariance and conditional independence give
\[
 \EE[f(\widehat O_1)f(\widehat O_2)]=\EE a_H^2,
 \qquad
 \EE[f(R_1)f(R_2)]=\EE d_H^2.
\]
The coupling bound \eqref{eq:switched-observation-comparison} and
$|f|\le1$ therefore imply
$|\EE a_H^2-\EE d_H^2|\le Cy^{-c}$.
With one base, switch invariance gives
$\EE[f(\widehat O_1)d_{\widehat H}]=\EE(a_Hd_H)$.
Replacing the multiplier by $d_H$ costs $Cy^{-c}$ by
\eqref{eq:multiplier-stability}; replacing $f(\widehat O_1)$ by
$f(R_1)$ costs the same by
\eqref{eq:switched-observation-comparison}.  Finally,
$\EE[f(R_1)d_H]=\EE d_H^2$, by conditioning on $H$.  Thus
\[
                   |\EE(a_Hd_H)-\EE d_H^2|\le Cy^{-c}.
\]
Combining these two comparisons yields
\begin{equation}\label{eq:empirical-test-variance}
                         \EE|a_H-d_H|^2\le Cy^{-c}.
\end{equation}
In particular the comparison is within each random empirical law,
not just after ensemble averaging.

For completeness, pass from tests to transportation as follows.
Embed the compact coordinate space in a fixed Euclidean cube and
use a tent partition of unity on a mesh $\delta=y^{-\epsilon}$.
There are $O(\delta^{-k})$ functions for a fixed dimension $k$, each
bounded by one and with Lipschitz constant $O(\delta^{-1})$.
Applying \eqref{eq:empirical-test-variance} to rescaled functions
bounds the expected sum of their mass discrepancies by
$C\delta^{-k-1}y^{-c/2}$.  Match the common mass in each tent and
couple the unmatched mass arbitrarily.  Since each tent has diameter
$O(\delta)$ and the whole space has bounded diameter, the resulting
squared transportation cost is at most a constant times $\delta^2$
plus that sum of discrepancies.  Choose $\epsilon$ small and apply
Markov's inequality.  This gives \eqref{eq:empirical-recursion}, with
a smaller positive exponent, for the cavity, full-root, and augmented
laws.  The proof of Proposition~\ref{prop:empirical-recursion} is complete.

\subsection{Descending the imaginary part}

The empirical recursion now supplies the hypotheses of the tree
stability theorem, except for its smallness condition.  We establish
that condition at a fixed positive height and retain it while descending.

\begin{proof}[Proof of Proposition~\ref{prop:local-law}]
At $\eta=h$, the cavity resolvent at a vertex with a sufficiently large
fixed tree neighborhood is close to the clean-tree cavity value,
first by polynomial approximation of the resolvent and then by the
norm perturbation bound $Cw/h^2$.  Equivalently one may approximate
the disordered resolvent locally and use its bounded potential
perturbation.  The polynomial degree is fixed after $h$ and the desired
accuracy are chosen; the Hamiltonians have uniformly bounded spectra.
The fraction of vertices failing this fixed tree test has expectation
$O(n^{-1})$.  Write $\delta_*$ for the fixed $L^2$ smallness threshold in
Proposition~\ref{prop:tree-stability}.  As permitted by its parameter
order, take $h$ small and then $w_0/h^2$ small enough that the starting
disk norm is below $\delta_*/2$, and take the target tree norm below
$\delta_*/4$.  Polynomial approximation with a sufficiently small
fixed error then places $P_H$ inside the threshold with failure
probability $O(n^{-1})$.

Fix $r\in(0,1)$ sufficiently close to one.  For any diagonal Green
function of a fixed graph or cut graph, Schwarz--Pick applied to its
map from $\HH$ to $\DD$ implies a uniform bound on its change between
$e+i\eta$ and $e+ir\eta$ tending to zero as $r\uparrow1$.
Choose $r$ so that this change is smaller than the gap between the
improved smallness bound and the threshold.  Starting at $h$, descend
through $h,rh,r^2h,\ldots$, shortening the last step to reach the
prescribed $\eta$.  At each height apply
Proposition~\ref{prop:empirical-recursion} for $j=1,\ell$ and then
Proposition~\ref{prop:tree-stability}.  Its conclusion places $P_H$
within $C(n\eta)^{-c}$ of $Q_\zeta$, hence back well inside the
threshold when $y_0$ is large.  This closes the induction.

The failure probabilities sum geometrically from the smallest height:
if $\eta_*$ is the terminal height, then
\[
             \sum_{\text{heights }\eta_k\ge\eta_*}
                       C(n\eta_k)^{-c}\le C'(n\eta_*)^{-c}.
\]
The initial $O(n^{-1})$ failure fits this bound after decreasing $c$.
We obtain \eqref{eq:distributional-local-law} for $P_H$.
For the full-root law, use the full-root version of
Proposition~\ref{prop:empirical-recursion} with one internal level.
Couple the cavity inputs to iid tree inputs optimally and apply the
sibling estimate \eqref{eq:tree-sibling-bound}, with $d$ children.
Its coefficients involve inverse margins of the other tree inputs,
whose required fixed moments are bounded.  Independence of the coupled
input pairs therefore bounds the output $L^2$ error by a constant
times the input $W_2$ error.  This proves the full statement.
\end{proof}

\subsection{Environment, trace, and microscopic integrability}

The disk metric was convenient for stability because it is bounded.
For later spectral and eigenvector arguments we need consequences in
half-plane coordinates, where large values must be controlled.  The
intercept $B_v$ in \eqref{eq:intercept} provides that control because
it is independent of the potential at the observed site.

\begin{corollary}[Consequences of the local law]
\label{cor:local-consequences}
Fix $E$ in the interior of $J$ and $s\in\RR$.  The following statements
hold.
\begin{enumerate}
\item At $\zeta=E+(s+iy)/n$, in the order $n\to\infty$ and then
$y\to\infty$, the diagonal $G_{vv}(\zeta)$ at a uniform vertex,
jointly with every fixed ordered neighborhood of potentials around
$v$, converges to the tree diagonal $m_o(E+i0)$ with the corresponding
tree environment.  For any fixed number of independently sampled
vertices the joint assertion holds with independent tree environments,
and their off-diagonal Green entries tend to zero in probability.
The diagonal assertion also holds along every sequence
$n\eta\to\infty$, $(e,\eta)\to(E,0)$.
\item There are $c,C>0$ such that, for every sufficiently large fixed $y$,
\begin{equation}\label{eq:trace-slow-power}
 \limsup_{n\to\infty}\PP\left\{
 \left|\Im\frac1n\Tr G\left(E+\frac{s+iy}{n}\right)
             -\pi\rho_{d,w}(E)\right|>y^{-c}\right\}\le Cy^{-c}.
\end{equation}
The constants are independent of the fixed real number $s$.
\item For each fixed $s\in\RR$ and $y>0$, the sampled diagonal imaginary
parts are uniformly integrable in the following limiting sense:
\begin{equation}\label{eq:microscopic-ui}
 \lim_{T\to\infty}\limsup_{n\to\infty}
 \EE_{H,v}\left[
 \Im G_{vv}\left(E+\frac{s+iy}{n}\right)
 \mathbf1_{\{\Im G_{vv}(E+(s+iy)/n)>T\}}\right]=0.
\end{equation}
\end{enumerate}
The double-limit assertion in the first item means that, for bounded
continuous tests, the limsup of the discrepancy as $n\to\infty$
tends to zero as $y\to\infty$.
\end{corollary}

\begin{proof}
For the first item, take the augmented recursion of
Proposition~\ref{prop:empirical-recursion} with internal levels covering
the prescribed neighborhood.  By Proposition~\ref{prop:local-law},
its input law tends to the tree cavity law.  Its internal iid potential
coordinates are independent of these inputs by construction.
The tree recursion with these data is exactly the joint tree law of
the diagonal and the observed neighborhood.  For any fixed additional
depth this passage only needs weak convergence: all limiting input
values lie in the open disk almost surely, where the finite recursion
is continuous.  Thus no higher inverse-margin order is required as the
observed radius grows.  Proposition~\ref{prop:tree-stability} supplies
the common-environment boundary continuity.  Conversion back to the
half-plane is valid since the limiting tree value is finite almost
surely.  The empirical-law conclusion gives product limits for
independently sampled vertices, and their mutual Green entries vanish
by \eqref{eq:sampled-ward}.  These arguments apply along every parameter
sequence stated in the corollary.  A diagonal choice of sequences gives
the stated double-limit formulation.  Letting the finite radius
increase identifies the limit with the measurable function
$m_o(E+i0)$ of its complete tree environment.

We next obtain a useful quantitative bound on small intercepts.
For a half-plane value $g$ and $x=X(g)$, the comparison used in
\eqref{eq:tree-hyperbolic-comparison} gives
\[
 \frac{\Im g}{|g|^2}\ge c_J(1-|x|).
\]
Couple $P_H^0$ with the tree law using
\eqref{eq:distributional-local-law}.  The bounded diameter of the disk
also bounds the expectation of its squared transportation error by
a negative power of $y$.  Markov's inequality for this coupling and
the tree inverse-margin moments show, for a sufficiently small
$\epsilon>0$, that
\begin{equation}\label{eq:small-intercept-large-height}
 \PP_{H,v}\{B_v(\zeta)<y^{-\epsilon}\}\le Cy^{-c}.
\end{equation}
For example, if the tree disk margin is at least $2\delta$ and the
coupling distance is below $\delta$, the finite margin is at least
$\delta$; take $\delta$ to be a sufficiently small power of $y^{-1}$.
This proves the bound directly from the two respective tail estimates.

The intercept formula gives $\Im G_{vv}\le B_v^{-1}$.
Since $B_v$ is independent of $V_v$, conditional spectral averaging
implies, for every $T>0$,
\begin{equation}\label{eq:intercept-tail-mean}
 \EE\bigl[\Im G_{vv}\mathbf1_{\{\Im G_{vv}>T\}}\bigr]
 \le \EE\bigl[\Im G_{vv}\mathbf1_{\{B_v<T^{-1}\}}\bigr]
 \le C_w\PP\{B_v<T^{-1}\}.
\end{equation}
This estimate applies after averaging over the uniform vertex as well.
Thus truncating the imaginary part at $T=y^\epsilon$ loses only a
negative-power amount in expectation, by
\eqref{eq:small-intercept-large-height}.  Tree tails have the same
property by their moment bounds.

To compare truncated imaginary parts through the disk local law,
first introduce a disk cutoff with margin $y^{-\epsilon'}$.
On the remaining compact subdisk, the function $x\mapsto\Im X^{-1}(x)$
has a Lipschitz constant bounded by a fixed power of $y^{\epsilon'}$.
Choose $\epsilon'$ small enough for the transportation estimate to
absorb this factor.  The mass removed by the cutoff has a negative-power
bound by the same coupling argument.  Choose the clipping exponent
$\epsilon$ still smaller, so that multiplying this mass by $y^\epsilon$
retains a negative-power bound.  Combining the comparisons and then
using Markov's inequality gives a negative-power bound for the
difference between the empirical imaginary trace and
$\EE\Im m_o(e+i\eta)$.  In the microscopic substitution with fixed
$s,y$, continuity of the tree expectation gives
$\EE\Im m_o\to\pi\rho_{d,w}(E)$ by
Corollary~\ref{cor:tree-density}.  This proves
\eqref{eq:trace-slow-power}, after decreasing $c$.  All finite-graph
constants were uniform in $e\in J$, hence are independent of fixed $s$.

Finally fix a possibly small microscopic height $y>0$.
The quantity $B_v(\zeta)$ is the imaginary part of the Herglotz function
\[
 \zeta+\langle a_v,(H^{(v)}-\zeta)^{-1}a_v\rangle,
\]
where $H^{(v)}$ deletes the vertex and $a_v$ is its adjacency column.
For $Y\ge y$, Harnack's inequality at the same real part yields
\[
 B_v\left(E+\frac{s+iy}{n}\right)
 \ge\frac yY B_v\left(E+\frac{s+iY}{n}\right).
\]
Choose a sufficiently large fixed $Y$.  Once
$\delta<(y/Y)Y^{-\epsilon}$, the probability that the left intercept
is below $\delta$ is at most the small-intercept probability at height
$Y$.  Consequently
\[
 \lim_{\delta\downarrow0}\limsup_n
 \PP\left\{B_v\left(E+\frac{s+iy}{n}\right)<\delta\right\}
 \le CY^{-c}.
\]
Let $Y\to\infty$ and apply \eqref{eq:intercept-tail-mean} to obtain
\eqref{eq:microscopic-ui}.  This last conclusion will retain the
normalization of eigenvector marks when taking microscopic limits.
\end{proof}
